% Options for packages loaded elsewhere
\PassOptionsToPackage{unicode}{hyperref}
\PassOptionsToPackage{hyphens}{url}
\PassOptionsToPackage{space}{xeCJK}
%
\documentclass[
  ignorenonframetext,
  aspectratio=169,
]{beamer}
\title{バイオフィルムの空間反応拡散モデル（Heine HOBIC）}
\subtitle{FISH 深さプロファイルからの空間輸送パラメータ逆推定 ---
形式的取り扱い}
\author{西岡佳祐 --- NIFE / SFB TRR-298}
\date{2026-06-03}

\usepackage{pgfpages}
\setbeamertemplate{caption}[numbered]
\setbeamertemplate{caption label separator}{: }
\setbeamercolor{caption name}{fg=normal text.fg}
\beamertemplatenavigationsymbolsempty
% Prevent slide breaks in the middle of a paragraph
\widowpenalties 1 10000
\raggedbottom
\setbeamertemplate{part page}{
  \centering
  \begin{beamercolorbox}[sep=16pt,center]{part title}
    \usebeamerfont{part title}\insertpart\par
  \end{beamercolorbox}
}
\setbeamertemplate{section page}{
  \centering
  \begin{beamercolorbox}[sep=12pt,center]{part title}
    \usebeamerfont{section title}\insertsection\par
  \end{beamercolorbox}
}
\setbeamertemplate{subsection page}{
  \centering
  \begin{beamercolorbox}[sep=8pt,center]{part title}
    \usebeamerfont{subsection title}\insertsubsection\par
  \end{beamercolorbox}
}
\AtBeginPart{
  \frame{\partpage}
}
\AtBeginSection{
  \ifbibliography
  \else
    \frame{\sectionpage}
  \fi
}
\AtBeginSubsection{
  \frame{\subsectionpage}
}
\usepackage{amsmath,amssymb}
\usepackage{lmodern}
\usepackage{iftex}
\ifPDFTeX
  \usepackage[T1]{fontenc}
  \usepackage[utf8]{inputenc}
  \usepackage{textcomp} % provide euro and other symbols
\else % if luatex or xetex
  \usepackage{unicode-math}
  \defaultfontfeatures{Scale=MatchLowercase}
  \defaultfontfeatures[\rmfamily]{Ligatures=TeX,Scale=1}
  \setmainfont[]{Noto Serif CJK JP}
  \setmonofont[]{Noto Sans Mono CJK JP}
  \ifXeTeX
    \usepackage{xeCJK}
    \setCJKmainfont[]{Noto Serif CJK JP}
  \fi
  \ifLuaTeX
    \usepackage[]{luatexja-fontspec}
    \setmainjfont[]{Noto Serif CJK JP}
  \fi
\fi
\usetheme[]{Madrid}
\usecolortheme{whale}
\usefonttheme{serif} % use mainfont rather than sansfont for slide text
% Use upquote if available, for straight quotes in verbatim environments
\IfFileExists{upquote.sty}{\usepackage{upquote}}{}
\IfFileExists{microtype.sty}{% use microtype if available
  \usepackage[]{microtype}
  \UseMicrotypeSet[protrusion]{basicmath} % disable protrusion for tt fonts
}{}
\makeatletter
\@ifundefined{KOMAClassName}{% if non-KOMA class
  \IfFileExists{parskip.sty}{%
    \usepackage{parskip}
  }{% else
    \setlength{\parindent}{0pt}
    \setlength{\parskip}{6pt plus 2pt minus 1pt}}
}{% if KOMA class
  \KOMAoptions{parskip=half}}
\makeatother
\usepackage{xcolor}
\IfFileExists{xurl.sty}{\usepackage{xurl}}{} % add URL line breaks if available
\IfFileExists{bookmark.sty}{\usepackage{bookmark}}{\usepackage{hyperref}}
\hypersetup{
  pdftitle={バイオフィルムの空間反応拡散モデル（Heine HOBIC）},
  pdfauthor={西岡佳祐 --- NIFE / SFB TRR-298},
  hidelinks,
  pdfcreator={LaTeX via pandoc}}
\urlstyle{same} % disable monospaced font for URLs
\newif\ifbibliography
\usepackage{longtable,booktabs,array}
\usepackage{calc} % for calculating minipage widths
\usepackage{caption}
% Make caption package work with longtable
\makeatletter
\def\fnum@table{\tablename~\thetable}
\makeatother
\usepackage{graphicx}
\makeatletter
\def\maxwidth{\ifdim\Gin@nat@width>\linewidth\linewidth\else\Gin@nat@width\fi}
\def\maxheight{\ifdim\Gin@nat@height>\textheight\textheight\else\Gin@nat@height\fi}
\makeatother
% Scale images if necessary, so that they will not overflow the page
% margins by default, and it is still possible to overwrite the defaults
% using explicit options in \includegraphics[width, height, ...]{}
\setkeys{Gin}{width=\maxwidth,height=\maxheight,keepaspectratio}
% Set default figure placement to htbp
\makeatletter
\def\fps@figure{htbp}
\makeatother
\setlength{\emergencystretch}{3em} % prevent overfull lines
\providecommand{\tightlist}{%
  \setlength{\itemsep}{0pt}\setlength{\parskip}{0pt}}
\setcounter{secnumdepth}{-\maxdimen} % remove section numbering
\usepackage{amsmath,amssymb}
\newcommand{\sgn}{\operatorname{sgn}}
\newcommand{\relu}[1]{\left[#1\right]_{+}}
\ifLuaTeX
  \usepackage{selnolig}  % disable illegal ligatures
\fi

\begin{document}
\frame{\titlepage}

\begin{frame}{このデッキの位置づけ}
\protect\hypertarget{ux3053ux306eux30c7ux30c3ux30adux306eux4f4dux7f6eux3065ux3051}{}
本プロジェクト = 3 本柱 ＋ 空間拡張。データの流れ：

raw 16S → guild \(\varphi\) → gLV/Hamilton（＋符号 prior）→ LOO 検証 →
空間 PDE

デッキ構成（◀ = 本デッキ）：

\begin{itemize}
\tightlist
\item
  \textbf{Overview}（傘）--- 全体像と3本柱
\item
  \textbf{AGORA} --- 代謝 → 符号 prior（生態モデルの入力）
\item
  \textbf{Dieckow} --- in-vivo 相互作用の推定と検証（生態モデル本体）
\item
  \textbf{Network} --- 相互作用行列 \(A\) の構造解析
\item
  \textbf{Spatial-PDE} --- FISH 深さプロファイルの反応拡散 ◀
  \textbf{本デッキ}
\item
  \textbf{FISH pipeline} --- .lif → 5 種深さ組成
\end{itemize}
\end{frame}

\begin{frame}{問い：WHO に WHERE を足す}
\protect\hypertarget{ux554fux3044who-ux306b-where-ux3092ux8db3ux3059}{}
bulk な ODE フィットは \textbf{誰が誰と相互作用するか（WHO）} を与える。
FISH 深さプロファイルは \textbf{それがどこで起こるか（WHERE）}
を加える。

\begin{itemize}
\tightlist
\item
  既知：bulk gLV/Hamilton の相互作用行列 \(A\)、増殖ベクトル
  \(b\)（Dieckow フィット）。
\item
  未知：種ごとの基質深さ方向の\textbf{輸送}（拡散・移流）。
\item
  目標：反応項を bulk
  から固定したまま、\textbf{空間輸送パラメータのみ}を 深さ分解 FISH
  データから逆推定する。
\end{itemize}

\vspace{0.4em}

\textbf{主張：} dysbiosis は bulk
組成の変化ではなく、\textbf{空間的再編成}として現れる。 特に \emph{P.
gingivalis} が深部へ沈み、\emph{F. nucleatum} から自律化する。
\end{frame}

\begin{frame}[fragile]{データ：HOBIC フローチャンバー FISH（Heine
2025）}
\protect\hypertarget{ux30c7ux30fcux30bfhobic-ux30d5ux30edux30fcux30c1ux30e3ux30f3ux30d0ux30fc-fishheine-2025}{}
\textbf{11 個の \texttt{.lif} ファイル}、2 条件 × 5 日 × チタン基板：

\begin{itemize}
\tightlist
\item
  条件：\textbf{CH}（commensal, HOBIC22）と \textbf{DH}（dysbiotic,
  HOBIC24）。
\item
  時点：Day 1 / 6 / 10 / 15 / 21。
\item
  5 種：So（\emph{S.
  oralis}）、An（\emph{Actinomyces}）、Vd/Vp（\emph{Veillonella}）、
  Fn（\emph{F. nucleatum}）、Pg（\emph{P. gingivalis}）。
\item
  プールした FOV 数：CH \(=15/5/15/16/16\)、DH \(=7/3/3/2/2\)。
\end{itemize}

\vspace{0.3em}

4 チャンネル FISH を 5 種にデコード（\emph{F. nucleatum} \(=\) 青
\(\cap\) 赤の二重標識； FISH デッキ参照）。z 軸 voxel 強度 \(\to\)
深さ方向の組成 \(\varphi_i(z,t)\)。
\end{frame}

\begin{frame}{深さプロファイル：CH vs DH}
\protect\hypertarget{ux6df1ux3055ux30d7ux30edux30d5ux30a1ux30a4ux30ebch-vs-dh}{}
\includegraphics[width=\textwidth,height=0.68\textheight]{results/diffusion_fit/zprofiles_all_ti_overlay.png}

種ごとの \(\varphi_i(z)\) を CH（commensal）と
DH（dysbiotic）で重ね描き。 深さ \(z\) は基質（\(z=0\)）から
bulk（\(z=L\)）方向。
\end{frame}

\begin{frame}{深さごとの積層組成}
\protect\hypertarget{ux6df1ux3055ux3054ux3068ux306eux7a4dux5c64ux7d44ux6210}{}
\includegraphics[width=\textwidth,height=0.68\textheight]{results/diffusion_fit/zprofiles_all_ti_stacked.png}

各深さ \(z\) における \(\{\varphi_i(z)\}_i\)
の積層表示。\(\sum_i \varphi_i + \varphi_0 = 1\) （\(\varphi_0\)
は空隙／非占有）を各 \(z\) で満たす。
\end{frame}

\begin{frame}{支配方程式（反応拡散 PDE）}
\protect\hypertarget{ux652fux914dux65b9ux7a0bux5f0fux53cdux5fdcux62e1ux6563-pde}{}
深さ \(z\in[0,L]\)、組成 \(\varphi_i(z,t)\) に対する反応 --- 移流 ---
拡散方程式： \[\frac{\partial \varphi_i}{\partial t}
 = D_i\,\frac{\partial^2 \varphi_i}{\partial z^2}
 \;-\; u\,\frac{\partial \varphi_i}{\partial z}
 \;+\; R_i(\varphi).\]

\begin{itemize}
\tightlist
\item
  \(R_i(\varphi)\) は replicator / Hamilton 反応項。\(A,b\) は bulk
  フィットから \textbf{固定}。
  \[R_i(\varphi)=\varphi_i\Big[(A\varphi+b)_i - \varphi^{\!\top}(A\varphi+b)\Big]+\gamma\,\varphi_i.\]
\item
  \textbf{自由パラメータは拡散率 \(D_i\) と移流 \(u\) のみ}。
\item
  \(\gamma\) は Lagrange 項：各 \(z\) で
  \(\sum_i \varphi_i + \varphi_0 = 1\) を強制する。
\end{itemize}
\end{frame}

\begin{frame}{境界条件と数値解法}
\protect\hypertarget{ux5883ux754cux6761ux4ef6ux3068ux6570ux5024ux89e3ux6cd5}{}
\textbf{境界条件：}
\[\left.\frac{\partial \varphi_i}{\partial z}\right|_{z=0}=0\quad(\text{基質，no-flux}),\qquad
\varphi_i\big|_{z=L}=\varphi_{\text{bulk},i}\quad(\text{bulk，Dirichlet}).\]
\(\varphi_{\text{bulk},i}\) は Day 1 の中央値で固定。

\textbf{数値解法（Method of Lines）：} Lie
演算子分割で反応と輸送を交互に解く。
\[e^{\Delta t\,\mathcal{L}} \approx e^{\Delta t\,\mathcal{R}}\,e^{\Delta t\,\mathcal{T}} + \mathcal{O}(\Delta t).\]

\begin{itemize}
\tightlist
\item
  反応 \(\mathcal{R}\)：各 \(z\) ノードで陰的 Euler（Newton 反復、stiff
  対応）。
\item
  輸送 \(\mathcal{T}\)：陽的有限差分（中心差分 \(\partial_z^2\)、風上
  \(\partial_z\)）。
\end{itemize}
\end{frame}

\begin{frame}{逆問題}
\protect\hypertarget{ux9006ux554fux984c}{}
観測 \(\varphi^{\text{obs}}(z,t)\)
に対する輸送パラメータの最小二乗推定： \[\min_{D,\,u}\;\; \sum_t \sum_z
 \big\lVert \varphi^{\text{pred}}(z,t;D,u) - \varphi^{\text{obs}}(z,t)\big\rVert^2 .\]

\begin{itemize}
\tightlist
\item
  最適化は \textbf{L-BFGS-B}。stiff PDE を通した数値勾配（JAX
  前向きモデル）。
\item
  前向きモデルは上記 MoL ソルバ。\(A,b\) 固定なので未知は
  \(D\in\mathbb{R}^5\) と \(u\)。
\item
  無次元化：\textbf{Péclet 数} \(\mathrm{Pe} = uL/D_i\) が移流 vs
  拡散の優劣を表す。 \(\mathrm{Pe}\gg1\) で移流支配、\(\mathrm{Pe}\ll1\)
  で拡散支配。
\end{itemize}
\end{frame}

\begin{frame}{予測 vs 観測：深さプロファイル}
\protect\hypertarget{ux4e88ux6e2c-vs-ux89b3ux6e2cux6df1ux3055ux30d7ux30edux30d5ux30a1ux30a4ux30eb}{}
\begin{columns}[T]
\begin{column}{0.48\textwidth}
\includegraphics[width=\textwidth,height=0.62\textheight]{results/diffusion_fit/fit_CH.png}

CH（commensal）
\end{column}

\begin{column}{0.48\textwidth}
\includegraphics[width=\textwidth,height=0.62\textheight]{results/diffusion_fit/fit_DH.png}

DH（dysbiotic）
\end{column}
\end{columns}

予測（実線）と観測（点）の深さプロファイル。反応項固定・輸送のみフィット。
\end{frame}

\begin{frame}{フィットされたパラメータ（暫定）}
\protect\hypertarget{ux30d5ux30a3ux30c3ux30c8ux3055ux308cux305fux30d1ux30e9ux30e1ux30fcux30bfux66abux5b9a}{}
\begin{longtable}[]{@{}lcc@{}}
\toprule
種 & \(D\)（CH） & \(D\)（DH） \\
\midrule
\endhead
So & \(0.0123\) & \(0.0596\) \\
An & \(0.0121\) & \(0.0021\) \\
Vd/Vp & \(0.0042\) & \(2.1\times10^{-5}\) \\
Fn & \(0.0075\) & \(0.0019\) \\
Pg & \(0.0115\) & \(0.0083\) \\
\textbf{\(u\)（移流）} & \(0.0038\) & \(0.0060\) \\
\textbf{loss} & \(0.128\) & \(0.102\) \\
\bottomrule
\end{longtable}

\vspace{0.2em}

\textcolor{red}{**暫定値**：両条件とも `success=False`（optimiser 未収束）。}
ハイパーパラメータ sweep と高速収束設定を HPC で実行中。
\end{frame}

\begin{frame}{空間生態 知見 1：Pg は dysbiosis で深く沈む}
\protect\hypertarget{ux7a7aux9593ux751fux614b-ux77e5ux898b-1pg-ux306f-dysbiosis-ux3067ux6df1ux304fux6c88ux3080}{}
\includegraphics[width=\textwidth,height=0.58\textheight]{results/diffusion_fit/depth_niche.png}

\emph{P. gingivalis} の質量中心（center-of-mass）が、DH では Day 6 以降
\textbf{最大 \(+30\,\mu\mathrm{m}\)
深く}シフトする。深部嫌気ニッチへの沈降。
\end{frame}

\begin{frame}{空間生態 知見 2：Fn--Pg ブリッジは初期のみ}
\protect\hypertarget{ux7a7aux9593ux751fux614b-ux77e5ux898b-2fnpg-ux30d6ux30eaux30c3ux30b8ux306fux521dux671fux306eux307f}{}
\includegraphics[width=\textwidth,height=0.5\textheight]{results/diffusion_fit/fn_pg_coloc.png}

Manders \(M_1\)（Pg のうち Fn と共局在する割合）： DH Day 1 \(=0.76\) →
Day 6 以降 \(0.15\text{–}0.20\)。CH は \(0.33\text{–}0.49\)。
\(\Rightarrow\) Pg は Fn から \textbf{脱共役}し、自律的に振る舞う。
\end{frame}

\begin{frame}{その他：bulk は CH ≈ DH}
\protect\hypertarget{ux305dux306eux4ed6bulk-ux306f-ch-dh}{}
CH--DH の \textbf{Bray--Curtis divergence} \(\approx 0.2\)、Day 1
から\textbf{ほぼ一定}。

\begin{itemize}
\tightlist
\item
  差は接種（inoculum）で決まり、時間発展でほとんど拡大しない。
\item
  bulk 組成では CH \(\approx\) DH \(\Rightarrow\) dysbiosis は
  \textbf{空間的・時間的}であり、 bulk 組成のシフトではない。
\end{itemize}

\vspace{0.4em}

これは知見 1・2 と整合する：違いは「誰がどれだけいるか」ではなく
「どこにいて、何と組むか」に現れる。
\end{frame}

\begin{frame}{拡散率 vs 生態的中心性}
\protect\hypertarget{ux62e1ux6563ux7387-vs-ux751fux614bux7684ux4e2dux5fc3ux6027}{}
\includegraphics[width=\textwidth,height=0.52\textheight]{results/diffusion_fit/d_vs_centrality.png}

フィットした拡散率 \(D_i\) は gLV 相互作用強度（中心性）と
\textbf{負相関}： CH \(r=-0.40\)、DH \(r=-0.47\)（5 種、illustrative）。

\vspace{0.2em}

\(\Rightarrow\)
「中心的な種ほど動かない」。生態的ハブが空間的に固定される傾向。
\end{frame}

\begin{frame}{in-vitro ↔ in-vivo マッピング}
\protect\hypertarget{in-vitro-in-vivo-ux30deux30c3ux30d4ux30f3ux30b0}{}
\includegraphics[width=\textwidth,height=0.5\textheight]{results/diffusion_fit/hobic_vs_dieckow.png}

HOBIC の 5 種を Dieckow の 5 ギルドに対応づけ、\textbf{順位が一致}：
Spearman \(\rho = 0.70\)。

\begin{itemize}
\tightlist
\item
  in-vivo は \emph{S. oralis} 優占、定義済み inoculum は均等接種。
\item
  順位構造は保たれる \(\Rightarrow\) in-vitro モデルは in-vivo
  の縮約として妥当。
\end{itemize}
\end{frame}

\begin{frame}{結論}
\protect\hypertarget{ux7d50ux8ad6}{}
\begin{enumerate}
\tightlist
\item
  bulk ODE（WHO）に FISH 深さプロファイル（WHERE）を足し、反応項固定で
  \textbf{輸送パラメータ \(D_i,u\) のみ}を逆推定する枠組みを構築。
\item
  \textbf{dysbiosis \(=\) 空間的再編成}：\emph{P. gingivalis}
  が深部へ沈み（\(+30\,\mu\mathrm{m}\)）、 \emph{F. nucleatum}
  から自律化（Manders \(M_1\) が初期のみ高い）。
\item
  bulk は CH \(\approx\) DH（Bray--Curtis \(\approx0.2\) で一定）---
  違いは bulk でなく空間。
\item
  「中心的な種ほど動かない」（\(D\) と中心性が負相関）；in-vitro↔in-vivo
  の 順位は一致（\(\rho=0.70\)）。
\item
  \textcolor{red}{輸送パラメータは**暫定**（未収束）}。HPC で
  sweep・高速収束設定が収束中。
\end{enumerate}
\end{frame}

\end{document}
